% Dossier para centros de yoga — Mantra Yoga
% Compilar con: xelatex dossier-centros.tex
\documentclass[11pt,a4paper]{article}

\usepackage{fontspec}
\usepackage{geometry}
\usepackage{xcolor}
\usepackage{colortbl}
\usepackage{amssymb}
\usepackage{tabularx}
\usepackage{booktabs}
\usepackage{titlesec}
\usepackage{enumitem}
\usepackage{fancyhdr}
\usepackage{microtype}

\geometry{a4paper, top=2cm, bottom=1.5cm, left=2.2cm, right=2.2cm}

% --- Colores -------------------------------------------------------------
\definecolor{tinta}{HTML}{1A1A1A}
\definecolor{ocre}{HTML}{9C5A28}
\definecolor{gris}{HTML}{6B6B6B}
\definecolor{grisclaro}{HTML}{F2F0ED}

% --- Tipografía ----------------------------------------------------------
\setmainfont{Palatino}
\newfontfamily\titular{Helvetica Neue}[UprightFont={* Medium}, BoldFont={* Bold}]

\color{tinta}
\raggedright
\setlength{\parindent}{0pt}
\setlength{\parskip}{0.45em}

% --- Secciones -----------------------------------------------------------
\titleformat{\section}
  {\titular\bfseries\fontsize{12}{14}\selectfont\color{ocre}}
  {}{0pt}{}
\titlespacing*{\section}{0pt}{1.5em}{0.4em}

\setlist[itemize]{leftmargin=1.1em, itemsep=0.25em, topsep=0.3em, label=\textcolor{ocre}{\footnotesize$\blacksquare$}}

% --- Pie de página -------------------------------------------------------
\pagestyle{fancy}
\fancyhf{}
\renewcommand{\headrulewidth}{0pt}
\renewcommand{\footrulewidth}{0pt}
\fancyfoot[L]{\titular\fontsize{7.5}{9}\selectfont\color{gris}Centros de Bhakti yoga · Entidad No Lucrativa · CIF G-76660679 · bhaktiyoga.es}
\fancyfoot[R]{\titular\fontsize{7.5}{9}\selectfont\color{gris}\thepage/2}

% --- Caja destacada ------------------------------------------------------
\newcommand{\caja}[1]{%
  \par\medskip
  \noindent\colorbox{grisclaro}{%
    \begin{minipage}{\dimexpr\textwidth-2\fboxsep}
      \vspace{0.4em}#1\vspace{0.4em}
    \end{minipage}}%
  \par\medskip}

\begin{document}

% =========================== PÁGINA 1 ===================================

{\titular\bfseries\fontsize{30}{34}\selectfont MANTRA YOGA\par}
\vspace{0.2em}
{\titular\fontsize{13}{16}\selectfont\color{gris} Un taller de una hora para tus alumnos,\\ sin coste para el centro\par}

\vspace{0.7em}
{\color{ocre}\rule{\textwidth}{1.6pt}}

\caja{\textbf{En una línea.} Una sesión mensual para aprender a cantar bien los
mantras clásicos en sánscrito. Guiada, gratuita para el centro y para quien
asiste.}

\section{Qué gana tu centro}

\begin{itemize}
  \item \textbf{Llena una franja floja del horario} sin pagar a nadie.
  \item \textbf{Algo distinto para tus alumnos} sin que tengas que montarlo tú.
  \item \textbf{No compite con tus clases.} Nadie deja el vinyasa por esto. Es
        complementario y se nota en la clase de asana de la semana siguiente.
  \item \textbf{Gente nueva conoce tu centro.} Los que vengan de fuera son
        alumnos potenciales tuyos, no nuestros.
  \item \textbf{Contenido listo para tus redes:} cartel y textos ya hechos.
\end{itemize}

\section{Qué te pedimos}

La sala, una hora, una vez al mes. Y que lo cuentes a tus alumnos.

Nada más. No hay alquiler, no hay porcentaje, no hay exclusividad y no hay
permanencia. \textbf{No se vende ni se publicita nada.} Si después del primero
no te convence, no hay un segundo.

\section{Qué es el taller}

Sesenta minutos, en círculo, en el suelo.

\vspace{0.4em}
{\setlength{\extrarowheight}{3pt}
\begin{tabularx}{\textwidth}{@{}>{\titular\bfseries\color{ocre}}l@{\hspace{1.3em}}X@{}}
\toprule
4 min  & El mantra del mes: cuál es, de dónde sale, qué dice \\
7 min  & Postura, respiración y preparación de la voz \\
12 min & Pronunciación: sílaba a sílaba y línea a línea \\
11 min & Cantarlo: al unísono primero, luego cada uno a su paso \\
5 min  & Asentamiento en silencio \\
11 min & Cápsulas cortas de cultura védica, entre bloque y bloque \\
4 min  & Preguntas, cierre y práctica para casa \\
\bottomrule
\end{tabularx}}

\vspace{0.6em}
Cada mes se trabaja un tema distinto y de él salen los mantras: el sonido de
AUM, R\=amacandra, la invocaci\'on del \=I\'sopani\d{s}ad. Cada sesión se entiende sola, así
que quien llega nuevo no se pierde y quien repite no repite. No hace falta
experiencia, ni saber cantar, ni saber sánscrito.

\newpage

% =========================== PÁGINA 2 ===================================

\section{De dónde viene}

El mantra yoga no es un añadido moderno. Patanjali lo pone dentro del método en
los \textit{Yoga-sutras} 1.27 y 1.28: el sonido y su repetición como práctica.

Lo que se canta pertenece a la tradición bhakti de la India, en el linaje de
Chaitanya (siglo XVI). Se explica con claridad en la sesión y se contesta a
cualquier pregunta sobre ello. No se hace proselitismo, no se pide adhesión a
nada, no se vende nada.

\section{Quién lo imparte}

\textbf{Juan Manuel Ferrera} (Jagannatha Mishra Dasa). Practico mantra yoga
desde 1981: cuarenta y cinco años cantando a diario. He guiado recitación y
meditación con mantra en España y fuera de ella.

\textbf{Centros de Bhakti yoga}, entidad no lucrativa registrada,
CIF G-76660679. bhaktiyoga.es

\section{Logística}

\vspace{0.4em}
{\setlength{\extrarowheight}{3pt}
\begin{tabularx}{\textwidth}{@{}>{\titular\bfseries}l@{\hspace{1.6em}}X@{}}
\toprule
Duración & 60 min (90 si prefieres una sesión más larga) \\
Frecuencia & Una vez al mes, día fijo \\
Aforo & De 8 a 30 personas \\
Espacio & Tu sala habitual. Cojines o mantas \\
Aportamos & Instrumentos, malas de préstamo y hojas para los asistentes \\
Coste para el centro & Ninguno \\
Coste para el alumno & Ninguno. Hay donativo voluntario, sin insistencia \\
Montaje & Llegamos 30 minutos antes \\
\bottomrule
\end{tabularx}}

\section{Lo que suelen preguntarme}

\textbf{¿Y si no viene nadie?} Entonces no has perdido nada, porque no has
pagado nada. Yo tampoco. Probamos uno más o lo dejamos.

\textbf{¿Vas a captar a mis alumnos?} Al revés. No tengo centro al que
llevarlos, y la gente que venga de fuera conoce el tuyo. A quien quiera seguir
le digo que vuelva aquí el mes que viene.

\textbf{¿Por qué es gratis?} Somos una entidad no lucrativa y el objetivo es que
esta práctica llegue a gente. Hay un plato de donativo voluntario al final, sin
insistir. No hay nada más.

\textbf{Mis alumnos no van a querer cantar.} Es lo que dicen todos los centros y
luego cantan. Empiezo muy bajito y aviso de que quien no quiera cantar puede
escuchar.

\vspace{0.6em}
{\color{ocre}\rule{\textwidth}{1.6pt}}
\vspace{0.6em}

{\titular\bfseries\fontsize{14}{17}\selectfont\color{ocre} Cómo seguimos\par}

\vspace{0.3em}
No hace falta que decidas nada ahora ni que contestes a esto. Me pasaré un día
de estos por el centro a saludar y, si pillo buen momento, te lo cuento en cinco
minutos y decides con calma. Si prefieres que avisemos antes, dime tú el día.

\vspace{0.5em}
{\titular\bfseries\fontsize{13}{16}\selectfont info@bhaktiyoga.es \quad·\quad 687 35 76 60\par}

\vspace{0.3em}
{\fontsize{10}{12}\selectfont\color{gris}Juan Manuel Ferrera · Centros de Bhakti yoga · bhaktiyoga.es}

\end{document}
